\chapter{Discussion transversale finale}

\textbf{Comment le deep learning adapte-t-il ses architectures à la structure
des données --- tabulaire, image et séquentielle --- et pourquoi un même
paradigme d'apprentissage supervisé doit-il être décliné différemment selon la
géométrie, la dépendance locale, la temporalité et la représentation des
données ?}

\paragraph{Unité du paradigme.}
Les trois parties du projet partagent un même cadre : l'apprentissage
différentiable supervisé par minimisation d'une fonction de coût
$\mathcal{L}(\theta)$ via descente de gradient stochastique. Le graphe de
calcul, la rétropropagation et l'optimiseur restent identiques en substance.

\paragraph{Diversité des biais inductifs.}
Ce qui change d'une architecture à l'autre est le \emph{biais inductif} que
chacune encode :
\begin{itemize}
    \item le MLP suppose que les entrées sont un vecteur plat de
          $\mathbb{R}^d$, sans structure interne ;
    \item le CNN suppose une grille spatiale 2D avec invariance approximative
          par translation et hiérarchie locale ;
    \item le RNN/LSTM/GRU suppose un ordre séquentiel avec dépendances
          contextuelles causales.
\end{itemize}
Formellement, le MLP est invariant à toute permutation des dimensions d'entrée
(modulo réarrangement des poids), le CNN est équivariant aux translations
(sortie translatée si l'entrée est translatée), et le RNN est sensible à
l'ordre causal de la séquence.

\paragraph{Conséquences expérimentales observées.}
\begin{itemize}
    \item Sur les données tabulaires (11 variables, Depression Professional Dataset), le MLP
          atteint un F1 $> 0{,}89$ car les variables sont déjà des descripteurs
          informatifs.
    \item Sur les images ($32 \times 32$ RGB, CIFAR-10), le CNN 
          atteint près de 73\,\% d'accuracy car il exploite la localité
          spatiale et les hiérarchies de traits visuels.
    \item Sur les séquences textuelles (traductions type IMDb), les architectures récurrentes
          modélisent des dépendances temporelles que ni le MLP ni le CNN
          standard ne capturent.
\end{itemize}

\paragraph{Synthèse.}
Le deep learning n'est donc pas un bloc monolithique mais une famille de
méthodes dont la pertinence dépend de l'adéquation entre le biais inductif
du modèle et la géométrie des données. L'unité réside dans l'apprentissage
de représentations par optimisation différentiable ; la diversité réside dans
les contraintes architecturales choisies pour rendre cet apprentissage
structurellement efficient.
